% Generated draft: requires mathematical and source review.
\documentclass[11pt]{article}
\usepackage[margin=25mm]{geometry}
\usepackage{fontspec}
\setmainfont{Latin Modern Roman}
\usepackage{amsmath,amssymb,mathtools,unicode-math}
\setmathfont{Latin Modern Math}
\usepackage{enumitem,needspace,xurl,hyperref}
\setlength{\parindent}{0pt}
\setlength{\parskip}{5pt}
\setlength{\emergencystretch}{3em}
\newcommand{\N}{\mathbb{N}}
\newcommand{\Z}{\mathbb{Z}}
\newcommand{\Q}{\mathbb{Q}}
\newcommand{\R}{\mathbb{R}}
\newcommand{\C}{\mathbb{C}}
\newcommand{\F}{\mathbb{F}}
\newcommand{\A}{\mathbb{A}}
\newcommand{\PP}{\mathbb{P}}
\newcommand{\E}{\mathbb{E}}
\newcommand{\eps}{\varepsilon}
\newcommand{\dd}{\,\mathrm{d}}
\begin{document}
\section*{UnsolvedMath ID: 30001038}
\textbf{Source problem code:} catalogue entry
\subsection{Definitions and mathematical statement}
In the standard orientation-preserving convention, let $G=\operatorname{IET}([0,1))$ be the group of bijections $T:[0,1)\to[0,1)$ for which there is a finite partition into half-open intervals and $T$ is a translation on each interval. The images partition $[0,1)$, and the operation is composition. No bound common to all group elements is placed on the number of intervals. The source does not explicitly specify whether exchanges with flips are allowed; the translation-only definition identifies the usual convention. The first question asks whether free groups embed into $G$, in particular whether there exists an injective abstract homomorphism \[F_2=\langle a,b\rangle\longrightarrow G.\] Equivalently for this specified rank, are there $T_1,T_2\in G$ such that every nonempty reduced word in $T_1^{\pm1},T_2^{\pm1}$ is nonidentity? The broader phrase free groups gives no rank range, so no assertion about arbitrarily large cardinal ranks is silently added. The second question asks which real special orthogonal groups \[\operatorname{SO}(n)=\{Q\in M_n(\R):Q^{\mathsf T}Q=I_n,\ \det Q=1\}\] embed into $G$. Embedding is read algebraically as an injective group homomorphism; no continuous Lie-group action or topology on $G$ is required by the statement. The source gives no range of $n$. The supplied 2026-08-21 triage labels the record partially solved: it reports that connected nonabelian Lie groups, including $\operatorname{SO}(n)$ for $n\geq3$, cannot embed, while the existence of a nonabelian free subgroup remains open. The reported exclusion is not extended to the lower-dimensional abelian or trivial orthogonal cases.
\subsection{Short English statement}
Does the interval-exchange group contain a free subgroup of rank two, and which special orthogonal groups occur as abstract subgroups? The supplied assessment reports nonembedding for $\operatorname{SO}(n)$ with $n\geq3$ while leaving the free-subgroup question open.
\subsection{Sources}
\begin{itemize}\item Local UnsolvedMath registry, permanent ID 30001038, source problem catalogue entry, \emph{Free and Orthogonal Subgroups of Interval Exchange Groups}, proposed year 2008; no individual proposer supplied.\item Supplied citation: \emph{Geometric Group Theory, Hyperbolic Dynamics and Symplectic Geometry} (2009), Oberwolfach Reports, DOI 10.4171/owr/2008/42, pp.\ 2337--2406: \url{https://doi.org/10.4171/owr/2008/42}. Supplied literature label: Oberwolfach report.\item Supplied literature label: Dahmani--Fujiwara--Guirardel, \url{https://arxiv.org/abs/1101.5909}.\item Supplied literature label: Free subgroups of IET, \url{https://arxiv.org/abs/1007.3940}.\item Literature assessment date: 2026-08-21. The status and progress claims are supplied dated curation, not independently verified here.\end{itemize}
\end{document}
